\chapter{Variations of Hodge structure on a general curve}
\label{chap:hodge-results}

% archon:covers MR4736527GeometricLocalSystemsOnVeryGeneralCurvesAndIsomonodromy/HodgeResults.lean

This chapter records the Hodge-theoretic results of \S7 of Landesman--Litt that
close the reduction chain. We isolate three lemmas: a criterion for unitarity of
the monodromy of a parabolically semistable variation of Hodge structure
(\cref{lem:unitary}); the passage from an integral representation that is unitary
under every embedding to one with finite image
(\cref{lem:integral-and-unitary-implies-finite}); and the non-density of
finite-image representations in the character variety
(\cref{lem:finite-image-not-dense}).

\section{The proof of the theorem on isomonodromic deformations underlying a CVHS}

\begin{lemma}[Semistable VHS has unitary monodromy]
  \label{lem:unitary}
  \lean{LandesmanLitt.unitary}
  \uses{cor:unstable, def:complex-variation, def:parabolic-stability,
    def:unitary-monodromy, prop:basic-facts, def:deligne-canonical-extension}
  % SOURCE: [MR4736527 / arXiv:2202.00039], Lemma (lemma:unitary), \S7.1
  % (read from references/MR4736527-local-systems-isomonodromy.tex, lines 3295-3303)
  % SOURCE QUOTE: "Suppose $(C, x_1, \ldots, x_n)$ is an $n$-pointed hyperbolic
  % curve and $(E,
  % \nabla)$ is a flat vector bundle on $C\setminus\{ x_1, \ldots, x_n\}$
  % underlying a polarizable complex variation of Hodge structure. Let
  % $(\overline{E}_\star, \overline{\nabla})$ be the Deligne canonical extension
  % of $(E, \nabla)$ to a flat vector bundle on $C$ with regular singularities at
  % the $x_i$, with the parabolic structure associated to $\overline{\nabla}$.
  % If $\overline{E}_\star$ is semistable, then the representation of
  % $\pi_1(C\setminus\{x_1, \cdots, x_n\})$
  % associated to $(E, \nabla)$ is unitary."
  \textit{Source: Landesman--Litt (arXiv:2202.00039), \S7, lemma:unitary.}

  Let $(C, x_1, \ldots, x_n)$ be an $n$-pointed hyperbolic curve and let
  $(E, \nabla)$ be a flat vector bundle on $C \setminus \{x_1, \ldots, x_n\}$
  underlying a polarizable complex variation of Hodge structure. Let
  $(\overline{E}_\star, \overline{\nabla})$ be the Deligne canonical extension of
  $(E, \nabla)$ to a flat bundle on $C$ with regular singularities at the $x_i$,
  equipped with the parabolic structure associated to $\overline{\nabla}$. If the
  parabolic bundle $\overline{E}_\star$ is semistable, then the monodromy
  representation of $\pi_1(C \setminus \{x_1, \ldots, x_n\})$ associated to
  $(E, \nabla)$ is unitary.
\end{lemma}

% SOURCE QUOTE PROOF: "By \autoref{prop:basic-facts}(2), we may write
% $\mathbb{V}:=\ker(\nabla)$ as $$\mathbb{V}:=\bigoplus_i
% \mathbb{L}_i\otimes W_i$$ where the $\mathbb{L}_i$ each have irreducible
% monodromy and also underlie polarizable variations of Hodge structure,
% and the $W_i$ are constant complex Hodge structures.
% It suffices to show the representation associated to each $\mathbb L_i$ has
% unitary monodromy.
% We may therefore reduce to the case that $\mathbb{V} = \mathbb L_i$ and assume
% that $(E,\nabla)$ has irreducible monodromy.
% Let $i$ be maximal such that $F^i\overline{E}_\star$ is non-zero.
% Since $\overline{E}_\star$ is semistable, it follows from
% \autoref{corollary:unstable} that
% that the natural map
% \begin{align*}
%   F^i\overline{E}_\star \to
%   F^{i-1}\overline{E}_\star/F^i\overline{E}_\star\otimes \omega_C
% \end{align*}
% induced by the connection is zero, i.e.~the connection preserves
% $F^i\overline{E}_\star$. By irreducibility
% of the monodromy of $(E, \nabla)$, we must have that $F^i\overline{E}_\star$
% equals $\overline{E}_\star.$
% But in this case $({E}, \nabla)$ is unitary, as the monodromy preserves the
% polarization, a definite Hermitian form."
\begin{proof}
  \uses{cor:unstable, prop:basic-facts, def:complex-variation,
    def:parabolic-stability, def:unitary-monodromy}
  By \cref{prop:basic-facts}(2) we may decompose the local system
  $\bV := \ker(\nabla)$ as $\bV = \bigoplus_i \bL_i \otimes W_i$, where each
  $\bL_i$ has irreducible monodromy and underlies a polarizable variation of
  Hodge structure and each $W_i$ is a constant complex Hodge structure. It
  suffices to prove that the representation associated to each $\bL_i$ is unitary,
  so we may reduce to the case $\bV = \bL_i$ and assume $(E, \nabla)$ has
  irreducible monodromy.

  Let $i$ be maximal such that $F^i \overline{E}_\star \neq 0$. Since
  $\overline{E}_\star$ is parabolically semistable, \cref{cor:unstable} shows that
  the map
  \[
    F^i \overline{E}_\star \to
      \bigl(F^{i-1}\overline{E}_\star / F^i\overline{E}_\star\bigr) \otimes
      \omega_C
  \]
  induced by the connection vanishes; that is, the connection preserves
  $F^i \overline{E}_\star$. Indeed, were the Hodge filtration to have more than
  one nonzero part, \cref{cor:unstable} would force $\overline{E}_\star$ to be
  parabolically unstable, contradicting semistability. By irreducibility of the
  monodromy, the flat subbundle $F^i \overline{E}_\star$ must equal
  $\overline{E}_\star$, so the Hodge filtration has a single nonzero piece. In
  this case the polarization is a flat, definite Hermitian form preserved by the
  monodromy, so the monodromy factors through the corresponding unitary group;
  hence $(E, \nabla)$ is unitary.
\end{proof}

\section{The proof of the rank bound for very general VHS}

\begin{lemma}[Integral and unitary implies finite image]
  \label{lem:integral-and-unitary-implies-finite}
  \lean{LandesmanLitt.integral_and_unitary_implies_finite}
  \uses{def:OK-local-system, def:unitary-monodromy}
  % SOURCE: [MR4736527 / arXiv:2202.00039], Lemma
  % (lemma:integral-and-unitary-implies-finite), \S7.2
  % (read from references/MR4736527-local-systems-isomonodromy.tex, lines 3384-3394)
  % SOURCE QUOTE: "Suppose $\Gamma$ is a group,
  % $K$ is a number field, and $\rho$ is
  % a representation $$\rho: \Gamma\to \on{GL}_m(\mathscr{O}_K).$$
  % If for each embedding $\iota: K\hookrightarrow \mathbb{C}$ the
  % representation $\rho\otimes_{\mathscr{O}_K, \iota} \mathbb{C}$ is unitary, then
  % $\rho$ has finite image."
  \textit{Source: Landesman--Litt (arXiv:2202.00039), \S7,
    lemma:integral-and-unitary-implies-finite.}

  Let $\Gamma$ be a group, let $K$ be a number field with ring of integers
  $\sO_K$, and let $\rho : \Gamma \to \GL_m(\sO_K)$ be a representation. If for
  every embedding $\iota : K \hookrightarrow \C$ the conjugate representation
  $\rho \otimes_{\sO_K, \iota} \C$ is unitary, then $\rho$ has finite image.
\end{lemma}

% SOURCE QUOTE PROOF: "Indeed, for $\iota: K\hookrightarrow\mathbb{C}$ an
% embedding, let $$\rho_\iota:
% \Gamma\to \on{GL}_m(\mathbb{C})$$ be the corresponding representation
% $\rho\otimes_{\mathscr{O}_K, \iota} \mathbb{C}$. First,
% $$\prod_\iota \rho_\iota: \Gamma\to \prod_\iota \on{GL}_m(\mathbb{C})$$ has
% image contained in a compact set by the definition of being unitary.
% Moreover, the image of $$\mathscr{O}_K\hookrightarrow \prod_\iota \mathbb{C}$$
% is discrete, since the difference of any two distinct elements of
% $\mathscr{O}_K$ has norm at least $1$.
% Hence the image of $\prod_\iota \rho_\iota$ is discrete and has compact closure,
% and is therefore finite."
\begin{proof}
  \uses{def:OK-local-system, def:unitary-monodromy}
  For each embedding $\iota : K \hookrightarrow \C$, write
  $\rho_\iota : \Gamma \to \GL_m(\C)$ for the conjugate representation
  $\rho \otimes_{\sO_K, \iota} \C$. By hypothesis each $\rho_\iota$ is unitary, so
  the product
  \[
    \textstyle\prod_\iota \rho_\iota : \Gamma \to \prod_\iota \GL_m(\C)
  \]
  has image contained in a product of compact unitary groups, hence relatively
  compact image. On the other hand, the diagonal image of
  $\sO_K \hookrightarrow \prod_\iota \C$, taken over all embeddings $\iota$, is
  discrete: the difference of any two distinct elements of $\sO_K$ is a nonzero
  algebraic integer, whose product of absolute values over all embeddings (its
  norm) is at least $1$. Therefore the image of $\prod_\iota \rho_\iota$ is a
  discrete subgroup with relatively compact closure, and is consequently finite.
  Since $\rho$ factors through this product, $\rho(\Gamma)$ is finite.
\end{proof}

\begin{lemma}[Finite-image representations are not Zariski dense]
  \label{lem:finite-image-not-dense}
  \lean{LandesmanLitt.finite_image_not_dense}
  \uses{def:character-variety}
  % SOURCE: [MR4736527 / arXiv:2202.00039], Lemma (lemma:finite-image-not-dense),
  % \S7.3
  % (read from references/MR4736527-local-systems-isomonodromy.tex, lines 3534-3540)
  % SOURCE QUOTE: "Let $\Sigma_{g,n}$ be a topological $n$-punctured genus $g$
  % surface with
  % basepoint $p \in \Sigma_{g,n}$. For $r > 1$, the set of representations
  % $\rho: \pi_1(\Sigma_{g,n},p) \to
  % \on{GL}_r(\mathbb C)$ with finite image are not Zariski dense in the character
  % variety $\mathscr{M}_{B,r}(\Sigma_{g,n}, p)$."
  \textit{Source: Landesman--Litt (arXiv:2202.00039), \S7,
    lemma:finite-image-not-dense.}

  Let $\Sigma_{g,n}$ be a topological $n$-punctured genus $g$ surface with
  basepoint $p \in \Sigma_{g,n}$. For $r > 1$, the representations
  $\rho : \pi_1(\Sigma_{g,n}, p) \to \GL_r(\C)$ with finite image are not Zariski
  dense in the character variety $\sM_{B,r}(\Sigma_{g,n}, p)$; they lie in a
  proper Zariski-closed subset.
\end{lemma}

% SOURCE QUOTE PROOF: "For $s\in \mathbb{Z}_{>0}$, let $V_s\subset
% \mathscr{M}_{B,r}(\pi_1(\Sigma_{g,n},p))$ be the closed subvariety
% consisting of those representations $\rho$ such that $$\on{Tr} [\rho(x)^s,
% \rho(y)^s]=\on{Tr}(\on{id})$$ for all $x,y\in \pi_1(\Sigma_{g,n},p)$.
% By Jordan's theorem (\cite[p.~91]{jordan1878memoire} or \cite[Theorem
% 36.13]{curtis1966representation}) on finite subgroups of
% $\on{GL}_r(\mathbb{C})$,
% there is some constant $m_r$ such that for each finite subgroup $G\subset
% \on{GL}_r(\mathbb{C})$, there exists an abelian normal subgroup $H\subset G$ of
% index at most $m_r$. Hence $V_{(m_r!)}$ contains all representations with finite
% image. Thus it suffices to show $V_s$ is not all of
% $\mathscr{M}_{B,r}(\Sigma_{g,n})$ for any $s>0$.
% We now write $$\pi_1(\Sigma_{g,n},p)=\left\langle a_1, \cdots, a_g, b_1,\cdots,
% b_g, \gamma_1, \cdots, \gamma_n \mid \prod_{i=1}^g [a_i,b_i]\cdot \prod_{j=1}^n
% \gamma_j\right\rangle$$ for the standard presentation of the fundamental group.
% If $g\geq 2$, let $$\rho: \pi_1(\Sigma_{g,n},p)\to \on{GL}_2(\mathbb{C})$$ be
% the representation defined by $$a_1 \mapsto \begin{pmatrix} 1 & 1 \\ 0 & 1
% \end{pmatrix}$$
% $$a_2 \mapsto \begin{pmatrix} 1 & 0 \\ 1 & 1 \end{pmatrix}$$
% and sending all other generators to the identity.  If $g=1$, hyperbolicity of
% $\Sigma_{g,n}$ implies $n>0$. Then $\pi_1(\Sigma_{g,n},x)$ is free on $n+1$
% generators, and we may set $\rho$ to be a representation sending two of the
% generators to the matrices above, and all other generators to the identity.
% Then $\rho$ does not lie in $V_s$ for any $s>0$ by direct computation. This
% completes the case $r = 2$. For the case that $r > 2$, the representation
% $\rho\oplus \on{triv}^{\oplus r-2}$ lies outside of $V_s$ for all $s$, for the
% same reason."
\begin{proof}
  \uses{def:character-variety}
  By Jordan's theorem on finite linear groups (Jordan, \emph{M\'emoire sur les
  \'equations diff\'erentielles lin\'eaires \`a int\'egrale alg\'ebrique}, p.~91;
  or Curtis--Reiner, Theorem 36.13), there is a constant $m_r$, depending only on
  $r$, such that every finite subgroup $G \subset \GL_r(\C)$ contains an abelian
  normal subgroup $H \subset G$ of index at most $m_r$. We take this result as a
  black box.

  For $s \in \Z_{>0}$, let $V_s \subseteq \sM_{B,r}(\Sigma_{g,n}, p)$ be the
  Zariski-closed subvariety of representations $\rho$ satisfying
  \[
    \on{Tr}\,[\rho(x)^s, \rho(y)^s] = \on{Tr}(\id)
    \qquad \text{for all } x, y \in \pi_1(\Sigma_{g,n}, p),
  \]
  where $[\,\cdot\,,\,\cdot\,]$ is the group commutator. If $\rho$ has finite
  image $G$, then raising to the power $s = m_r!$ lands every element in the
  abelian normal subgroup $H$, so the commutators $[\rho(x)^s, \rho(y)^s]$ are
  trivial; hence $V_{(m_r!)}$ contains every finite-image representation. It thus
  suffices to exhibit, for each $s > 0$, a single representation lying outside
  $V_s$, which shows $V_s \subsetneq \sM_{B,r}(\Sigma_{g,n}, p)$ is a proper
  closed subset.

  Using the standard presentation
  \[
    \pi_1(\Sigma_{g,n}, p) = \Bigl\langle a_1, \ldots, a_g, b_1, \ldots, b_g,
      \gamma_1, \ldots, \gamma_n \;\Big|\;
      \textstyle\prod_{i=1}^g [a_i, b_i] \cdot \prod_{j=1}^n \gamma_j
    \Bigr\rangle,
  \]
  consider, when $g \ge 2$, the representation
  $\rho : \pi_1(\Sigma_{g,n}, p) \to \GL_2(\C)$ sending
  \[
    a_1 \mapsto \begin{pmatrix} 1 & 1 \\ 0 & 1 \end{pmatrix},
    \qquad
    a_2 \mapsto \begin{pmatrix} 1 & 0 \\ 1 & 1 \end{pmatrix},
  \]
  and all other generators to the identity. When $g = 1$, hyperbolicity forces
  $n > 0$, so $\pi_1(\Sigma_{g,n}, p)$ is free on $n + 1$ generators and we send
  two generators to the matrices above and the rest to the identity. A direct
  computation shows $\rho \notin V_s$ for any $s > 0$, settling the case $r = 2$.
  For $r > 2$, the representation $\rho \oplus \on{triv}^{\oplus (r-2)}$ lies
  outside every $V_s$ by the same computation. Hence the finite-image
  representations are contained in the proper closed subvariety $V_{(m_r!)}$ and
  are not Zariski dense.
\end{proof}
